% Compile with:  pdflatex example.tex   (or lualatex example.tex)
\documentclass[11pt]{article}
\usepackage[dotgrid]{messynotes}

\begin{document}

\notehead[9/15/26]{Cook--Levin}

\headu{Warmup: mapping reductions}

\pt{3-COLOR $\le_p$ 4-COLOR}
\pt{$G = (V,E) \longmapsto G' = \bigl(V \cup \{s\},\; E \cup \{(s,v) : v \in V\}\bigr)$}

\pt{PALINDROME $\le_p$ 3-COLOR}
\pt[2]{$x \longmapsto$ 3-clique if palindrome, 4-clique otherwise.}

\pt{3-COLOR $\le_p$ HALT}

\imp{\lead{Fact} 3-COLOR is NP-complete}

\lead{Corollary} if 3-COLOR $\le_p$ PALINDROME $\Rightarrow$ P${}={}$NP

\lead{Pf} to show P${}={}$NP, exhibit a poly-time algorithm for arbitrary
$L \in$ NP under the assumed reduction

\begin{ind}
  Fixing $L \in$ NP $\Rightarrow \exists$ poly-time $R_{L \to 3\text{COL}}$
  by \amber{\ding{72}}

  Given instance $x$ to $L$:

  \bul{run $y \leftarrow R_{L \to 3\text{COL}}(x)$}
  \bul{run $z \leftarrow R_{3\text{COL} \to \text{PAL}}(y)$ \aside{by assumption}}
  \bul{solve $z \in$ PALINDROME}
  \bul[1]{if yes, accept}
  \bul[1]{if no, reject \quad\tomb}
\end{ind}

\headu{Cook--Levin Theorem}

SAT is NP-complete.

\arr{\key{Karp} showed an incredibly diverse set of problems $\Pi$ with
     SAT $\le_p \Pi$.}

\gap

\headbar{SAT: Boolean satisfiability}

Boolean formula in conjunctive normal form

\arr{$n$ boolean variables $x_1,\dots,x_n$}
\pt{$m$ clauses $c_j = (\ell_{j_1} \vee \ell_{j_2} \vee \cdots \vee \ell_{j_k})$}
\arr[1]{$m = \mathrm{poly}(n)$}
\pt[3]{$\ell_i \in \{x_i, \neg x_i\}$}

\eq{\phi \;=\; \bigwedge_{j=1}^{m} c_j}

\lead{SAT} does there exist an assignment $a \in \{0,1\}^n$
s.t.\ $x = a \Rightarrow \phi(a) = 1$?

\lead{$k$-SAT} each clause has width at most $k$.

\begin{ex}[Example: 4-SAT $\le_p$ 3-SAT]
Given a 4-CNF formula $\phi$, construct a 3-CNF formula $\psi$ such that
$\phi$ is satisfiable $\iff$ $\psi$ is satisfiable.

\pt{For each $c_j = (\ell_1 \vee \ell_2 \vee \ell_3 \vee \ell_4)$ in $\phi$,}
\pt[2]{$c_{j,+} = (\ell_1 \vee \ell_2 \vee y_j)$}
\pt[2]{$c_{j,-} = (\ell_3 \vee \ell_4 \vee \neg y_j)$ \aside{for brand new $y_j$}}

\eq[2]{\psi \;=\; \bigwedge_{j=1}^m c_{j,+} \wedge c_{j,-} \qquad \tomb}
\end{ex}

\sep

\headbox{Loose ends}

\begin{cols}{2}
\subhead{Still unclear}
\qn{why does the $y_j$ trick preserve satisfiability in both directions?}
\qn{is the $k \to 3$ blowup really linear?}
\nextcol
\subhead{To do}
\yes{reread Sipser 7.4}
\no{write up the $k$-SAT induction}
\bang{pset due Friday}
\end{cols}

\begin{quo}[Sipser, Thm 7.37]
SAT is NP-complete. The proof simulates a nondeterministic polynomial-time
machine with a tableau of configurations.
\end{quo}

\clearpage
\notehead[ref]{Everything else in the kit}

\headu{Bullets}
\bul{\texttt{\char`\\bul} -- plain dot}
\dash{\texttt{\char`\\dash} -- en dash}
\arr{\texttt{\char`\\arr} -- the ``sub-point'' arrow}
\imp{\texttt{\char`\\imp} -- star, for things that matter}
\qn{\texttt{\char`\\qn} -- something I did not follow}
\bang{\texttt{\char`\\bang} -- warning to future me}
\yes{\texttt{\char`\\yes} \quad and \quad \texttt{\char`\\no}}
\no{each takes an optional depth: \texttt{\char`\\bul[2]\{...\}}}

\newpoint{\hmm}{\textcolor{mnviolet}{$\rightsquigarrow$}}
\hmm{roll your own with \texttt{\char`\\newpoint\{\char`\\hmm\}\{$\rightsquigarrow$\}}}

\headu{Headings}
\head{\texttt{\char`\\head} -- plain}
\headu{\texttt{\char`\\headu} -- ruled}
\headbox{\texttt{\char`\\headbox} -- boxed}
\headbar{\texttt{\char`\\headbar} -- filled bar}
\subhead{\texttt{\char`\\subhead} -- small}

\headu{Panels}

\begin{defn}
A language $L$ is \key{NP-hard} if every $L' \in$ NP reduces to it in
polynomial time.
\end{defn}

\begin{idea}
Every panel takes an optional label: \texttt{\char`\\begin\{idea\}[My label]}.
Pass \texttt{[]} to drop the label entirely.
\end{idea}

\begin{warn}
The reduction has to run in polynomial time \alert{in the size of the
original input}, not the size of its output.
\end{warn}

\begin{outline}[Plain outlined box]
\texttt{outline} and \texttt{shade} are the unlabelled generics.
\end{outline}

\begin{shade}[]
Make your own: \texttt{\char`\\newshadebox\{hmm\}\{mnviolet\}\{Hmm\}} or
\texttt{\char`\\newoutlinebox\{...\}}.
\end{shade}

\headu{Side by side}

\sbs[0.55]{%
  \subhead{Wider left}
  \texttt{\char`\\sbs[0.55]\{left\}\{right\}} splits two blocks, optionally
  unevenly.
}{%
  \subhead{Narrower right}
  \bul{bullets work inside}
  \bul{so do boxes and math}
}

\begin{cols}{3}
\subhead{one}
$\sum_{i} a_i$
\nextcol
\subhead{two}
$\bigcup_{j} B_j$
\nextcol
\subhead{three}
$\vdash \preceq \ltimes$
\end{cols}

\headu{Emphasis}

\pt[0]{\key{key term}, \alert{alert}, \hl{highlighted}, \aside{muted aside},
\ul{underlined}, \blue{blue} \teal{teal} \red{red} \amber{amber}
\violet{violet} \gray{gray}}

\sep

\aside{\texttt{\char`\\sep} draws that rule; \texttt{\char`\\gap} adds blank
vertical space; \texttt{\char`\\eq\{...\}} sets left-aligned display math at the
current indent.}

\end{document}
